%% toc: P1 8 8.3,8.4,8.6,8.7
\chapter{Reading the Equivalence Glyphs Across the Book}
\chaptermark{Equivalence Glyphs Across the Book}
\label{sec:p01-w02b-p01-013-equivalence-audit}
\index{equivalence sign!typed roles}\index{definitional sign}
\index{symbol registry!equivalence glyphs}

\paragraph{Purpose.}  The triple bar \(\equiv\), the double arrow
\(\Leftrightarrow\) and the definitional sign \(\defeq\) occur in many parts
of the book, and no single meaning fits all of them.  The registry of
\cref{sec:p01-symbol-registry} therefore records each use as a typed role
with a scope.  This chapter states the roles that the earlier conventions of
\cref{sec:p01-symbol-registry,sec:p01-w02-p01-025-signs} do not yet name,
gives the rule by which a reader recovers the role of any occurrence from its
operands, proves the two facts on which the new roles depend, and lists every
role of the three glyphs with the region of the book in which it is in force.

\paragraph{Inputs and conclusions.}  The input is one occurrence of
\(\equiv\) (with or without a subscript or superscript), of
\(\Leftrightarrow\), or of the definitional sign, together with the section in
which it stands.  The conclusion is the role of that occurrence: the kind of
its operands, the relation or connective it denotes, and whether it belongs to
the object language or to the metalanguage.

\section{Roles not covered by the earlier conventions}
\label{sec:p01-w02b-p01-013-new-roles}

The editorial defaults and the source-local roles recorded in
\cref{def:p01-w02-p01-025-equivalence-roles} remain in force.  The following
four conventions complete them.

\begin{hbdefinition}{def:p01-w02b-p01-013-congruence-predicate}{Presburger congruence predicate}
Let \(m\geq1\) and \(0\leq r<m\) be natural numbers.  In the language
\(\{0,1,+,<\}\) of \cref{def:p05-presburger}, write \(\underline r\) for
\(0\) when \(r=0\) and for the term \(1+\cdots+1\) with \(r\) summands when
\(r\geq1\), and write \(m\cdot y\) for the term \(y+\cdots+y\) with \(m\)
summands.  The \emph{congruence predicate} \(x\equiv r\pmod m\) is the unary
predicate symbol \(C_{m,r}\) added to the language together with its explicit
definition
\[
  C_{m,r}(x)\ \leftrightarrow\ \exists y\,\bigl(x=\underline r+m\cdot y\bigr),
\]
where \(y\) is a variable distinct from \(x\).  In this role \(m\) and \(r\)
are fixed numerals of the metalanguage and the operand \(x\) is a term of the
object language; the sign \(\equiv\) is part of the name of a predicate
symbol.  The role is in force in \cref{def:p05-presburger},
\cref{thm:p05-presburger-qe} and \cref{def:p05-presburger-example}.
\index{congruence predicate}\index{Presburger arithmetic!congruence predicates}
\end{hbdefinition}

The metalinguistic congruence \(t\equiv a\pmod n\) of item~6 of
\cref{def:p01-w02-p01-025-equivalence-roles} relates two numbers; the
predicate above is a symbol of the object language whose interpretation in
\(\mathbb N\) is that relation with the second argument fixed
(\cref{thm:p01-w02b-p01-013-congruence-definable}).  The same metalinguistic
congruence, with modulus \(2\), appears in \cref{sec:p09-presentation-n-c-012}.

\begin{hbconvention}{def:p01-w02b-p01-013-identically-equal}{Identically equal}
Let \(D\) be a nonempty set, let \(h:D\to\mathbb C\) be a function and let
\(c\in\mathbb C\).  Then \(h\equiv c\), read ``\(h\) is identically
\(c\)'', means \(h(y)=c\) for every \(y\in D\).  The left operand is a
function and the right operand a constant; the domain is the one declared for
\(h\).  The role is metalinguistic.  It is in force for \(h\equiv1/2\) in
\cref{thm:p15-conditional-examples}, for \(\varphi\equiv1\) in
\cref{thm:p15-characteristic-second-moment}, for \(\varphi\equiv1\) and
\(f\equiv0\) in \cref{sec:p15-characteristic-function-criteria-and-moments}
and its annex \cref{sec:annex-p15-characteristic-function-criteria-and-moments},
and for \(\Psi_m\equiv m/\delta\) in
\cref{sec:annex-p15-bayes-factor-computation}.
\index{identically equal}
\end{hbconvention}

\begin{hbconvention}{def:p01-w02b-p01-013-statement-iff}{Double arrow between assertions}
When both operands of \(\Leftrightarrow\) are assertions of the metalanguage,
such as \(x\in A\), \(\mathsf D\derives\sigma\) or
\((M,\in)\entails\varphi(\vec a)\), the sign is the biconditional of the
metalanguage: the displayed assertions are both true or both false.  When both
operands are formulas, \(\varphi\Leftrightarrow\psi\) is the editorial formula
equivalence of \cref{sec:p01-symbol-registry}.  The operands decide between
the two readings; a formula is never an assertion, since an assertion about a
formula needs a relation such as \(\entails\) or \(\derives\).  The
assertion reading is in force throughout the book; it is used, for example, in
\cref{sec:p04-complexity}, \cref{sec:p04-turing-degrees},
\cref{sec:p04-frontiers-computability}, \cref{def:p06-finite-character},
\cref{thm:p06-suslin}, \cref{def:p06-absoluteness} and
\cref{def:p01-w02-p01-022-setting}.
\index{biconditional!metalinguistic}
\end{hbconvention}

\begin{hbconvention}{def:p01-w02b-p01-013-mention}{Displayed register entries}
In the register tables of \cref{sec:p01-symbol-registry},
\cref{sec:p01-w02-p01-024-notation} and
\cref{sec:p14-symbols-terminology-and-glossary}, and in the typed-role
convention of \cref{sec:p01-w02-p01-025-overview}, an equivalence glyph is
displayed together with its registered role and scope, in the same table row
or the same sentence.  Such a display names the glyph and asserts no relation.
Its role is the one printed beside it.
\end{hbconvention}

\section{Recovering the role of an occurrence}
\label{sec:p01-w02b-p01-013-procedure}

\begin{hbmethod}{def:p01-w02b-p01-013-role-reading}{Reading an equivalence glyph}
\emph{Input:} an occurrence of \(\equiv\), \(\Leftrightarrow\) or the
definitional sign, and its section.
\begin{enumerate}
  \item If the sign is the definitional sign (printed \(\mathrel{:\equiv}\),
  or \(\mathrel{:=}\) inside the restoration scopes described in item~3 of
  \cref{def:p01-w02-p01-025-equivalence-roles}), the left expression is
  defined to abbreviate the right one.  Stop.
  \item If the sign carries a subscript or superscript, read it as one
  symbol and look it up in the registry under that full form.  The subscript
  is part of a name (\(\equiv_\alpha\), \(\equiv_m\), \(\equiv_{\mathrm j}\),
  \(\equiv_{\mathrm W}\)) or an index of a declared family
  (\(\equiv_E^n\), \(\equiv_k\), \(\equiv_Q\)); it is never a modulus.  Stop.
  \item If a modulus is displayed by \(\pmod{\cdot}\), the role is
  congruence: metalinguistic between numbers (item~6 of
  \cref{def:p01-w02-p01-025-equivalence-roles}) or the object-language
  predicate of \cref{def:p01-w02b-p01-013-congruence-predicate} inside the
  Presburger chapter.  Stop.
  \item If both operands are structures, the role is elementary equivalence.
  Stop.
  \item If the left operand is a function and the right a constant, the role
  is \cref{def:p01-w02b-p01-013-identically-equal}.  Stop.
  \item If the sign is \(\Leftrightarrow\), apply
  \cref{def:p01-w02b-p01-013-statement-iff}, unless the section declares it
  as a connective of its object language (as the internal higher-order
  language of \cref{sec:p10-n-d-023-source-03} does).  Stop.
  \item Otherwise the operands are formulas, terms of an object language or
  concepts.  If the section's object language lists \(\equiv\) among its
  connectives, it is that connective.  Otherwise it is the equivalence of
  the section's declared semantics: the two sides hold at the same points of
  every model of that semantics.
\end{enumerate}
\emph{Side conditions.}  Step~7 needs the semantics or object language to be
declared before the occurrence; \cref{tab:p01-w02b-p01-013-roles} names the
declaration for every region of the book.  Steps~4 and~5 read the operand
types from the declarations of the operands, not from their letters.
\emph{Output:} the role of the occurrence.
\end{hbmethod}

The steps are tried in order, but at most one of them applies to a given
occurrence, because their operand conditions exclude one another: numbers
with a displayed modulus, structures, a function against a constant,
assertions and formulas are five different kinds of object.  Where a section
uses \(\equiv\) both for the equivalence of its semantics and to introduce a
derived operator, as the temporal presentations of
\cref{sec:p09-presentation-n-c-012} do with \(FA\equiv\top\mathbin{\mathsf
U}A\), the two readings agree: once the operator is introduced by the line,
the two sides are true at the same points of every model.

\section{Facts behind the new roles}
\label{sec:p01-w02b-p01-013-facts}

\begin{hblemma}{thm:p01-w02b-p01-013-congruence-definable}{proved=annex:p01-w02b-p01-013-congruence-definable}{Congruence predicates are definitional}
Let \(m\geq1\) and \(0\leq r<m\), and let \(C_{m,r}\) be the predicate of
\cref{def:p01-w02b-p01-013-congruence-predicate}.
\begin{enumerate}
  \item For every \(n\in\mathbb N\),
  \(\mathbb N\entails\exists y\,(x=\underline r+m\cdot y)\) under the
  assignment \(x\mapsto n\) iff \(m\) divides the integer \(n-r\).
  \item For every \(n\in\mathbb N\) there is exactly one \(r<m\) for which
  the condition of part~1 holds.
  \item Let \(\mathbb N^{+}\) be the expansion of \((\mathbb
  N;0,1,+,<)\) in which each \(C_{m,r}\) is interpreted by
  \(\{n\in\mathbb N:m\text{ divides }n-r\}\).  For a formula \(\chi\) of the
  expanded language let \(\chi^{\circ}\) replace each atom \(C_{m,r}(t)\) by
  \(\exists y\,(t=\underline r+m\cdot y)\), where \(y\) is a variable that does
  not occur in \(\chi\).  Then for every assignment \(s\),
  \(\mathbb N^{+}\entails\chi[s]\) iff \(\mathbb N\entails\chi^{\circ}[s]\).
\end{enumerate}
\end{hblemma}

Part~3 is what makes the predicates a definitional expansion: adding them
changes no truth value of a formula of the original language and makes no
new relation definable, so \cref{thm:p05-presburger-qe} is a statement about
the original additive structure.

\begin{hblemma}{thm:p01-w02b-p01-013-identically-equal}{proved=annex:p01-w02b-p01-013-identically-equal}{Identical equality is pointwise}
Let \(D\) be a nonempty set, \(h:D\to\mathbb C\) and \(c\in\mathbb C\).
\begin{enumerate}
  \item \(h\equiv c\) iff \(h\) is the constant function \(\kappa_c\) with
  \(\kappa_c(y)=c\) for all \(y\in D\).
  \item Let \((D,\mathcal D,\mu)\) be a measure space, let \(y_0\in D\) with
  \(\{y_0\}\in\mathcal D\) and \(\mu\{y_0\}=0\), and let \(c'\neq c\).  If
  \(h\equiv c\) and \(h'\) agrees with \(h\) except that \(h'(y_0)=c'\), then
  \(h'=h\) \(\mu\)-almost everywhere and \(h'\not\equiv c\).
\end{enumerate}
\end{hblemma}

Part~2 is the situation of item~3 of \cref{thm:p15-conditional-examples}:
the version \(h\equiv1/2\) and the version altered at \(0\) are both versions
of the conditional mean, but only the first is identically \(1/2\).

\begin{hbexample}{def:p01-w02b-p01-013-example}{Three lines with the same glyph}
(1)~In \cref{def:p05-presburger-example}, \(x\equiv1\pmod3\) carries a
displayed modulus, so step~3 of \cref{def:p01-w02b-p01-013-role-reading}
applies inside the Presburger chapter: it is the atom \(C_{3,1}(x)\).  By
part~3 of \cref{thm:p01-w02b-p01-013-congruence-definable} it may be replaced
by \(\exists y\,(x=1+(y+y+y))\) without changing truth in \(\mathbb N\); the
example's formula \(\exists y\,(x=3y+1)\) is this formula with the sum
\(y+y+y\) written \(3y\) and the order of the summands exchanged, which does
not change its value.
(2)~In \cref{thm:p15-characteristic-second-moment}, \(\varphi\equiv1\) has a
characteristic function on the left and a constant on the right, so step~5
applies: \(\varphi(t)=1\) for every real \(t\).
(3)~In \cref{thm:p15-conditional-examples}, \(h\equiv1/2\) likewise means
\(h(y)=1/2\) for every \(y\in[0,1]\), not merely for almost every \(y\).
\end{hbexample}

\begin{hbwarning}{def:p01-w02b-p01-013-failure}{Readings that change the statement}
Reading \(h\equiv1/2\) as equality almost everywhere weakens item~3 of
\cref{thm:p15-conditional-examples}: by part~2 of
\cref{thm:p01-w02b-p01-013-identically-equal} the altered version is equal to
\(h\) almost everywhere, yet the theorem uses the difference between the two
versions at the null point \(0\).  Reading \(x\equiv r\pmod m\) in the
Presburger chapter as a metalinguistic statement about a number loses the
fact that quantifier elimination needs these atoms in the object language.
Reading any of these lines as elementary equivalence or as formula equivalence
is a type error, since their operands are neither two structures nor two
formulas.  The procedure has a limit: it reads operand types from
declarations, so an occurrence whose operands are not declared before it has
no determined role, and the registry gives none.
\end{hbwarning}

\section{All roles of the three glyphs}
\label{sec:p01-w02b-p01-013-table}

\Cref{tab:p01-w02b-p01-013-roles} lists every role in which the three glyphs
occur in the reading text, the kind of operands, and where the role is in
force.  An occurrence outside the listed regions with a given operand kind
does not have that role.  The definitional sign has a single role throughout,
printed \(\mathrel{:\equiv}\) in editorial text and \(\mathrel{:=}\) inside the
restoration scopes.

\begingroup\small
\begin{longtable}{@{}p{0.22\textwidth}p{0.26\textwidth}p{0.44\textwidth}@{}}
\caption{Roles of \(\equiv\), \(\Leftrightarrow\) and the definitional sign,
with their operand kinds and regions.}
\label{tab:p01-w02b-p01-013-roles}\\
\toprule
Form & Role and operands & Region \\
\midrule
\endfirsthead
\toprule
Form & Role and operands & Region \\
\midrule
\endhead
\(a\mathrel{:\equiv}b\), \(a\mathrel{:=}b\) & definitional sign; expressions &
everywhere; the second printing only in restoration scopes \\
\(\varphi\Leftrightarrow\psi\) & editorial formula equivalence; formulas &
everywhere \\
\(P\Leftrightarrow Q\) & metalinguistic biconditional; assertions &
everywhere (\cref{def:p01-w02b-p01-013-statement-iff}) \\
\(s\Leftrightarrow t\) & defined biconditional of the internal higher-order
language; terms of type \(\Omega\) & \cref{sec:p10-n-d-023-source-03},
\cref{annex:p10.mlm.8} \\
\(\Leftrightarrow\) & defined formula equivalence; formulas of the
propositional language of the prime-ideal chapter & the propositional
sections of \cref{sec:p01-boolean-prime-ideal-and-ultrafilters} and their
annex \\
\(\mathcal M\equiv\mathcal N\) & elementary equivalence; structures &
everywhere \\
\(\varphi\equiv\psi\) & CPL semantic equivalence; CPL formulas &
\cref{sec:prop-syntax-semantics-01}, \cref{sec:p02-w02-p02-007-equivalence-laws},
\cref{sec:annex-p02-w02-p02-007-equivalence-laws},
\cref{sec:p02-w02-p02-010-duality-arbitrary},
\cref{sec:annex-p02-w02-p02-010-duality-arbitrary},
\cref{sec:annex-p02-w02-p02-004-equivalent-substitutions},
\cref{annex:p2.1.8}, \cref{sec:p02-w02-p02-019-sat-procedures},
\cref{sec:annex-p02-w02-p02-019-sat-procedures},
\cref{sec:p02-w02-p02-016-finite-input-scope},
\cref{sec:annex-p02-w02-p02-016-finite-input-scope} \\
\(\varphi\equiv\psi\) & first-order logical equivalence; first-order formulas &
item~1 of \cref{def:p01-w02-p01-025-equivalence-roles};
\cref{sec:annex-p01-w02-p01-025-equality} \\
\(\alpha\equiv\beta\) & equivalence joining lines of a propositional
calculation & \cref{sec:writing-foundations-05} \\
\(\varphi\equiv\psi\) & schematic entry relative to a declared class &
item~5 of \cref{def:p01-w02-p01-025-equivalence-roles} \\
\(A\equiv B\) & classical logical equivalence; first-order formulas &
\cref{sec:p16-deductive-nomological-explanation},
\cref{sec:annex-p16-deductive-nomological-explanation} \\
\(U\equiv V\) & logical equivalence of sentences of a finite language &
\cref{sec:p15-systematic-power}, \cref{sec:annex-p15-systematic-power} \\
\(\equiv\) bare, in chains & abbreviation of \(\equiv_\Sigma\) &
\cref{sec:p12-small-clause-standard}, \cref{annex:p12-small-clause-standard} \\
\(A\equiv B\) & equivalence of the declared semantics (temporal, CTL, modal
\(\mu\), dynamic, separation, predicate-transformer, linear-order and
finite-tree formulas), also introducing defined operators &
\cref{sec:p09-logic-in-computer-science},
\cref{sec:p09-presentation-n-c-012}, \cref{sec:p09-presentation-n-c-013},
\cref{sec:p09-presentation-n-c-014}, \cref{sec:p09-presentation-n-c-016},
\cref{sec:p09-n-c-017-biere-ectl-certificates},
\cref{sec:p09-presentation-n-c-018}, \cref{sec:p09-presentation-n-c-019},
\cref{sec:p09-presentation-n-c-020}, \cref{sec:p09-presentation-n-c-021},
\cref{sec:p09-stavi-completeness}, \cref{sec:annex-p09-stavi-completeness},
\cref{annex:p9.trees.1} \\
\(C\equiv D\) & concept equivalence; description-logic concepts &
\cref{sec:p09-presentation-n-c-011} \\
\(\equiv\) & object-language biconditional connective & the causal theories
of \cref{sec:p16-nonmonotonic-causal-theories} and
\cref{sec:p16-cplus-action-language}; the erotetic languages of
\cref{sec:p18-erotetic-patterns-and-composition} and
\cref{sec:p18-erotetic-patterns-proofs}; the native search procedures of
Part~18; Reichenbach's logic in \cref{sec:p15-reichenbach-probability-logic};
Church's type theory in \cref{sec:p02-w02-p02-040-church} and
\cref{sec:annex-p02-w02-p02-040}; the notation-variant table of
\cref{sec:p01-notation-variants} \\
\(P_t\equiv\ldots\) & predicate definition &
\cref{def:p15-pr-learning-native-diagonal} and the native Gold predicates \\
\(\equiv\) & judgmental equality & the HoTT source theory of Part~7 \\
\(\equiv\) & local Boolean congruence; filter agreement of index sets &
the prime-ideal source dictionary of Part~1; the topological Boolean index
algebra of Part~3 \\
\(\equiv\) & native logical identity atom & the Delhomme presentation of
Part~3 \\
\(s\equiv\alpha\) & satisfaction of a formula by a state & the preferential
and ranked models of Part~8 \\
\(c_i\equiv c_j\) & equality pattern; constant indices &
\cref{sec:annex-p02-w02-p02-035-decision-procedures}; \(G_\equiv\) in
\cref{sec:p02-w02-p02-035-bsr} \\
\(x\equiv r\pmod m\) & Presburger congruence predicate; object-language term &
\cref{def:p05-presburger} through \cref{def:p05-presburger-example} \\
\(t\equiv a\pmod n\) & congruence of numbers & \cref{sec:decision-hierarchy-02},
\cref{sec:p09-presentation-n-c-012} \\
\(h\equiv c\) & identically equal; function and constant &
\cref{def:p01-w02b-p01-013-identically-equal} \\
\(\equiv_\alpha\) & alpha-equivalence; formulas &
\cref{sec:fol-syntax-semantics-04} and the later source chapters,
\cref{sec:rule-sheets-01}, \cref{sec:annex-p01-w02-p01-003-language-change},
\cref{sec:annex-p01-w02-p01-025-equality},
\cref{sec:annex-p02-w02-p02-037-craig-beth},
\cref{sec:p07-classical-cut-elimination},
\cref{sec:p07-equality-lk-metatheory},
\cref{sec:p07-intuitionistic-cut-elimination}, \cref{annex:p7.eqlk.alpha}
through \cref{annex:p7.eqlk.nd} \\
\(\equiv_\alpha\) & pointed back-and-forth equivalence at a countable
ordinal; tuples in structures & \cref{def:p03-morley-approximants} and its
results, \cref{annex:p3.morley.scott-sentence} \\
\(\equiv_m\) & many-one equivalence; sets of numbers &
\cref{sec:decision-hierarchy-03}, \cref{sec:p04-frontiers-computability} \\
\(\equiv_m\) & agreement on sentences of quantifier rank at most \(m\) &
\cref{annex:p3.abstract.3} \\
\(\equiv_F\) & agreement modulo a filter & \cref{sec:p03-fv-reduced-composition} \\
\(\equiv_{\mathrm W},\equiv_{\mathrm{sW}}\) & Weihrauch and strong Weihrauch
equivalence; problems & \cref{sec:p04-weihrauch-choice} and its annex \\
\(\equiv_{\mathrm j}\) & judgmental equality; terms &
\cref{sec:p07-higher-inductive-types}, \cref{sec:p07-hott-circle-encode-decode},
\cref{sec:p07-hott-identity-systems-and-flattening} and their annexes \\
\(\Gamma\vdash a\equiv b:A\) & judgmental equality of Martin-L\"of type
theory & \cref{sec:p07-lambda-calculi-and-type-theories} \\
\(\equiv_P\) & two-valued equivalence over arbitrary atoms &
\cref{sec:p02-w02-p02-016-finite-input-scope},
\cref{sec:annex-p02-w02-p02-016-finite-input-scope} \\
\(\equiv_L,\equiv_S\) & provable equivalence in a logic or system &
\cref{sec:p10-algebraic-and-categorical-logic} and the algebraic chapters of
Part~10 \\
\(\equiv_L\) & syntactic context congruence of a tree language &
the finite-tree algebra of Part~9 \\
\(\equiv_E^n\) & equational congruence in a finite context &
\cref{sec:p10-lawvere-theories} and its annex \\
\(\equiv_Q\) and instances & congruence generated by equations &
\cref{sec:p12-rewriting}, \cref{sec:p12-fair-completion} and their annexes \\
\(\equiv_\Sigma\) & pointwise first-order semantic equivalence &
\cref{sec:p12-small-clause-standard} \\
\(\equiv_n\), \(\equiv_k\), \(\equiv_{\Delta}\), \(\equiv_h\) & game
equivalence of orders; agreement of trees to a depth; filtration equivalence
of states; abstraction kernels & the temporal presentations of Part~9,
\cref{sec:p09-presentation-n-c-018}, \cref{sec:p09-presentation-n-c-021},
\cref{sec:p09-presentation-n-c-017} \\
\(\equiv^{\mathcal K},\equiv_s^{\mathcal K}\) & state and model-validity
equivalence over a class \(\mathcal K\) &
\cref{sec:p09-presentation-n-c-016-emerson-halpern} \\
\(\equiv_{\mathrm{InqB}}\) & mutual support entailment &
the inquisitive chapter of Part~8 \\
\(\equiv\), \(\Leftrightarrow\) displayed & register entry &
\cref{def:p01-w02b-p01-013-mention} \\
\bottomrule
\end{longtable}
\endgroup

\paragraph{Relationships.}  The roles that relate formulas are linked by the
passages of \cref{thm:p01-w02-p01-025-equiv-bridges}: first-order equivalence
restricts to CPL equivalence on propositional formulas, alpha-equivalence
implies first-order equivalence, and equivalent sentences keep their truth
value between elementarily equivalent structures.  The Presburger predicate
and the metalinguistic congruence are linked by part~1 of
\cref{thm:p01-w02b-p01-013-congruence-definable}.  Identical equality implies
equality almost everywhere and not conversely, by
\cref{thm:p01-w02b-p01-013-identically-equal}.  No other passage between two
rows of \cref{tab:p01-w02b-p01-013-roles} holds in general, because the
operands are of different kinds.

\paragraph{Sources.}  The Presburger predicate is the one used in
\cref{def:p05-presburger} for the elimination theorem of
\cite{Presburger1929}; the identically-equal reading is the one used in
\cref{thm:p15-conditional-examples} and
\cref{thm:p15-characteristic-second-moment}.  Both lemmas of this chapter are
proved in full in the annex and depend on no source beyond the division
algorithm for natural numbers and the definition of satisfaction.
